%% toc: P3 40 40.6
\section{Proofs for almost-sure encoding and infinite spectra}
\label{sec:annex-p03-w03-p03-069-infinite-spectra}

Throughout, \(G(n,p)\) is the random graph of \cref{def:p03-w03-p03-065-gnp}:
its edge indicators, one for each unordered pair of distinct vertices of
\([n]\), are independent and each equals \(1\) with probability \(p\).  For a
vertex set \(Z\), the \emph{pairs meeting \(Z\)} are the unordered pairs with
at least one end in \(Z\).  ``For large \(n\)'' means for all \(n\) beyond a
bound depending only on the fixed parameters named in the statement.

\begin{annexitem}{annex:p03-w03-p03-069-k3k-safe}{thm:p03-w03-p03-069-k3k-safe}{\(K_{3,k}\) as a rooted graph}
\begin{proof}
Write the vertex set of \(K_{3,k}\) as \(P\cup Q\) with \(|P|=3\),
\(|Q|=k\), the edges being all pairs \(\{x,y\}\) with \(x\in P\), \(y\in Q\).
There are no roots, so every vertex is a nonroot and every edge is counted:
\(v=3+k\) and \(e=3k\).  The rooted graph is sparse exactly when
\(3+k-3k\alpha>0\), that is, when
\(\alpha<(3+k)/(3k)=\tfrac13+\tfrac1k\).

Safe implies sparse, by taking \(S=V(H)\) in the definition of safe
(\cref{def:p03-w03-p03-065-dense-sparse}).  Conversely, suppose
\(\alpha<\tfrac13+\tfrac1k\), and let \(S\) be a nonempty subset of
\(P\cup Q\), with \(a=|S\cap P|\) and \(b=|S\cap Q|\).  The restriction of
\(K_{3,k}\) to \(S\) has exactly \(ab\) edges, one for each pair of an element
of \(S\cap P\) and an element of \(S\cap Q\).  If \(a=0\) or \(b=0\), then
\(|S|-ab\,\alpha=|S|>0\).  If \(a,b\geq1\), then \(a\leq3\) and \(b\leq k\)
give \(\tfrac1a+\tfrac1b\geq\tfrac13+\tfrac1k>\alpha\), so
\[
  (a+b)-ab\,\alpha=ab\Bigl(\tfrac1a+\tfrac1b-\alpha\Bigr)>0 .
\]
So every subextension \((\emptyset,K_{3,k}|_S)\) is sparse, and
\((\emptyset,K_{3,k})\) is safe.
\end{proof}
\end{annexitem}

\begin{annexitem}{annex:p03-w03-p03-069-k3k-threshold}{thm:p03-w03-p03-069-k3k-threshold}{Threshold for a \(K_{3,k}\)}
\begin{proof}
Let \(\Omega\) be the set of ordered pairs \(c=(P_c,Q_c)\) of disjoint subsets
of \([n]\) with \(|P_c|=3\) and \(|Q_c|=k\).  Let \(E(c)\) be the set of the
\(3k\) pairs \(\{x,y\}\) with \(x\in P_c\), \(y\in Q_c\), let \(B_c\) be the
event that every pair in \(E(c)\) is an edge, and let
\(X=\sum_{c\in\Omega}\mathbf 1_{B_c}\).  A \(K_{3,k}\) exists exactly when
\(X>0\).  Each \(B_c\) has probability \(p^{3k}\), and
\(|\Omega|=\binom n3\binom{n-3}k\leq n^{3+k}\).  For \(n\geq2(k+3)\), using
\(\binom mj\geq(m-j+1)^j/j!\),
\[
  |\Omega|\geq\frac{(n-k-3)^{3+k}}{3!\,k!}\geq d_k\,n^{3+k},
  \qquad d_k=\frac{1}{2^{3+k}\,3!\,k!}.
\]
Note that \((3+k)/(3k)=\tfrac13+\tfrac1k\), so
\(n^{3+k}p^{3k}=\bigl(p\,n^{1/3+1/k}\bigr)^{3k}\).

\emph{Below the threshold.}  If \(p\ll n^{-1/3-1/k}\), then
\[
  \mathbb E[X]\leq n^{3+k}p^{3k}=(p\,n^{1/3+1/k})^{3k}\to0,
\]
and
\(\Pr[X>0]\leq\mathbb E[X]\to0\) by the first moment bound of
\cref{thm:p03-w03-p03-065-moments}.

\emph{Above the threshold.}  Suppose \(p\gg n^{-1/3-1/k}\) and put
\(\omega=p\,n^{1/3+1/k}\), so \(\omega\to\infty\); also \(p\leq1\).  For
\(n\geq2(k+3)\), \(\mathbb E[X]\geq d_k\,\omega^{3k}\to\infty\).  If
\(E(c)\cap E(c')=\emptyset\) then \(B_c\) and \(B_{c'}\) are determined by
disjoint sets of independent edge indicators and are independent.  Let \(D\)
be the set of ordered pairs \((c,c')\), \(c\neq c'\), with
\(E(c)\cap E(c')\neq\emptyset\).  By \cref{thm:p03-w03-p03-065-moments},
\[
  \operatorname{Var}[X]\leq\mathbb E[X]+\sum_{(c,c')\in D}\Pr[B_c\wedge B_{c'}] .
\]
Fix \((c,c')\in D\), let \(W=(P_c\cup Q_c)\cap(P_{c'}\cup Q_{c'})\), and let
\(a=|W\cap P_c|\), \(b=|W\cap Q_c|\).  Every pair in \(E(c)\cap E(c')\) is a
pair of \(E(c)\) with both ends in \(W\), and there are exactly \(ab\) such
pairs; since \(E(c)\cap E(c')\neq\emptyset\), \(a\geq1\) and \(b\geq1\).
Hence \(|E(c)\cup E(c')|\geq6k-ab\) and, as \(p\leq1\),
\(\Pr[B_c\wedge B_{c'}]=p^{|E(c)\cup E(c')|}\leq p^{6k-ab}\).

For fixed \(c\) there are at most \(2^{3+k}\) choices of \(W\subseteq
P_c\cup Q_c\).  For fixed \(c\) and \(W\), every \(c'\) with
\(W\subseteq P_{c'}\cup Q_{c'}\) is fixed by the \(3+k-|W|\) vertices of
\(P_{c'}\cup Q_{c'}\) outside \(W\), at most \(n^{3+k-|W|}\) choices, and by
the choice of which three of its \(3+k\) vertices form \(P_{c'}\), at most
\(2^{3+k}\) choices.  Grouping the pairs of \(D\) by \(c\) and by the type
\((a,b)\) of \(W\), with \(1\leq a\leq3\) and \(1\leq b\leq k\),
\[
  \sum_{(c,c')\in D}\Pr[B_c\wedge B_{c'}]
  \leq|\Omega|\sum_{a=1}^{3}\sum_{b=1}^{k}4^{3+k}\,n^{3+k-a-b}\,p^{6k-ab}.
\]
Dividing by \(\mathbb E[X]^2=|\Omega|^2p^{6k}\) and using
\(|\Omega|\geq d_kn^{3+k}\),
\[
  \frac{\operatorname{Var}[X]}{\mathbb E[X]^2}
  \leq\frac1{\mathbb E[X]}
  +\sum_{a=1}^{3}\sum_{b=1}^{k}\frac{4^{3+k}}{d_k}\cdot
  \frac{1}{n^{a+b}p^{ab}} .
\]
For \(1\leq a\leq3\), \(1\leq b\leq k\), we have
\(ab(\tfrac13+\tfrac1k)\leq ab(\tfrac1a+\tfrac1b)=a+b\), so
\[
  n^{a+b}p^{ab}=\omega^{ab}\,n^{(a+b)-ab(1/3+1/k)}\geq\omega^{ab}\to\infty .
\]
The sum has \(3k\) terms, each tending to \(0\), and
\(1/\mathbb E[X]\to0\).  So \(\Pr[X=0]\leq\operatorname{Var}[X]/\mathbb
E[X]^2\to0\) by \cref{thm:p03-w03-p03-065-moments}, and a \(K_{3,k}\) exists
almost surely.
\end{proof}
\end{annexitem}

\begin{annexitem}{annex:p03-w03-p03-069-max-common}{thm:p03-w03-p03-069-max-common}{The largest common neighbourhood}
\begin{proof}
First, for every graph \(G\) with at least three vertices and every
\(j\geq1\), \(\nu(G)\geq j\) exactly when \(G\) contains a \(K_{3,j}\).
If \(|N(x_1,x_2,x_3)|\geq j\) for distinct \(x_1,x_2,x_3\), any \(j\)
elements of \(N(x_1,x_2,x_3)\) are distinct from the \(x_i\) by
irreflexivity and form a \(K_{3,j}\) with them.  Conversely the \(y\)'s of a
\(K_{3,j}\) lie in \(N(x_1,x_2,x_3)\).

Let \(\alpha\in I_k\) and \(p=n^{-\alpha}\).  Since
\(\alpha<\tfrac13+\tfrac1k\),
\(p/n^{-1/3-1/k}=n^{1/3+1/k-\alpha}\to\infty\), so by
\cref{thm:p03-w03-p03-069-k3k-threshold} a \(K_{3,k}\) exists almost surely
and \(\nu\geq k\) almost surely.  Since
\(\alpha>\tfrac13+\tfrac1{k+1}\),
\(p/n^{-1/3-1/(k+1)}=n^{1/3+1/(k+1)-\alpha}\to0\), so a \(K_{3,k+1}\) exists
almost never and \(\nu\leq k\) almost surely.  The intersection of two
events holding almost surely holds almost surely
(\cref{def:p03-w03-p03-065-almost-sure}).

\emph{The claim of \cref{def:p03-w03-p03-069-distinctness}.}  Let
\(\alpha\in(0,1)\), \(p=n^{-\alpha}\), and let \(X\) be the degree of the
vertex \(1\) in \(G(n,p)\).  Then \(X\) is binomial with \(n-1\) trials and
success probability \(p\), with mean \(\mu=(n-1)p\geq n^{1-\alpha}/2\) for
\(n\geq2\), so \(\mu\to\infty\), and variance \((n-1)p(1-p)\leq\mu\).  By
Chebyshev's inequality,
\(\Pr[X\leq\mu/2]\leq\Pr[|X-\mu|\geq\mu/2]\leq4\mu/\mu^2=4/\mu\to0\).  So
almost surely the vertex \(1\) has more than \((n-1)p/2\) neighbours, and
\(|N(1,1,1)|\), the number of its neighbours, tends to infinity in
probability.
\end{proof}
\end{annexitem}

\begin{annexitem}{annex:p03-w03-p03-069-even-bigger}{thm:p03-w03-p03-069-even-bigger}{What the second-order formulas say}
\begin{proof}
\emph{EVEN.}  Suppose \(R\) is irreflexive and symmetric and every \(x\in S\)
has exactly one \(y\in S\) with \(R(x,y)\); call it \(\pi(x)\).  Then
\(\pi(x)\neq x\) by irreflexivity.  Also \(R(\pi(x),x)\) by symmetry, and
\(x\in S\), so uniqueness of the partner of \(\pi(x)\) gives
\(\pi(\pi(x))=x\).  Thus \(\pi\) is an involution of \(S\) without fixed
points, its orbits \(\{x,\pi(x)\}\) have exactly two elements and partition
\(S\), and \(|S|\) is even.  Conversely, if \(|S|\) is even, partition \(S\)
into two-element sets and let \(R\) consist of the ordered pairs \((a,b)\) and
\((b,a)\) for each such set \(\{a,b\}\).  This \(R\) is irreflexive and
symmetric, and each element of \(S\) has exactly one partner in \(S\).

\emph{BIGGER.}  Suppose \(R\) witnesses \(\operatorname{BIGGER}(S,T)\).  For
each \(y\in T-S\) choose \(f(y)\in S-T\) with \(R(y,f(y))\).  If
\(f(y_1)=f(y_2)=x\) for distinct \(y_1,y_2\in T-S\), then \(R(y_1,x)\) and
\(R(y_2,x)\), which the second clause forbids.  So \(f\) is injective and
\(|T-S|\leq|S-T|\).  Conversely, if \(|T-S|\leq|S-T|\), take an injection
\(f\colon T-S\to S-T\) and let \(R\) consist of the pairs \((y,f(y))\) and
\((f(y),y)\) for \(y\in T-S\).  It is symmetric, and irreflexive because
\(y\notin S\) and \(f(y)\in S\).  Every \(y\in T-S\) has \(R(y,f(y))\).  If
\(R(y_1,x)\) and \(R(y_2,x)\) with \(y_1,y_2\in T-S\) and \(x\in S-T\), then
\(x=f(y_1)=f(y_2)\), because the only pairs of \(R\) with first coordinate in
\(T-S\) are the pairs \((y,f(y))\); injectivity gives \(y_1=y_2\).  Finally
\(|T|=|T\cap S|+|T-S|\) and \(|S|=|S\cap T|+|S-T|\), so
\(|T-S|\leq|S-T|\) exactly when \(|T|\leq|S|\).

\emph{\(A^*\).}  By the two equivalences, a graph \(G\) with at least three
vertices satisfies \(A^*\) exactly when some distinct triple \(\bar x\) has
\(|N(\bar x)|\) even and \(|N(\bar z)|\leq|N(\bar x)|\) for every distinct
triple \(\bar z\); that is, exactly when some distinct triple attains the
maximum \(\nu(G)\) and that maximum is even.  The maximum is attained, since
there are finitely many triples, so \(G\entails A^*\) exactly when \(\nu(G)\)
is even.
\end{proof}
\end{annexitem}

\begin{annexitem}{annex:p03-w03-p03-069-astar-spectrum}{thm:p03-w03-p03-069-astar-spectrum}{The limiting behaviour of \(A^*\)}
\begin{proof}
Let \(\alpha\in I_k\) and \(p=n^{-\alpha}\).  For \(n\geq3\),
\(G(n,p)\entails A^*\) exactly when \(\nu(G(n,p))\) is even, by
\cref{thm:p03-w03-p03-069-even-bigger}.  By
\cref{thm:p03-w03-p03-069-max-common}, \(\Pr[\nu(G(n,p))=k]\to1\).  If \(k\) is
even, \(\Pr[A^*]\geq\Pr[\nu=k]\to1\); if \(k\) is odd,
\(\Pr[A^*]\leq\Pr[\nu\neq k]\to0\).  In particular \(f_{A^*}\) is defined at
every irrational \(\alpha\in(1/3,1)\) other than the points \(1/3+1/k\), which
are rational, and \(f_{A^*}(\alpha)\) is \(1\) or \(0\) according as the
\(k\) with \(\alpha\in I_k\) is even or odd.

\emph{The points \(1/3+1/k\), \(k\geq2\).}  Let \(q=\tfrac13+\tfrac1k\), so
\(q\leq\tfrac56<1\).  The interval \(I_k\) has right end point \(q\), and
\(I_{k-1}\) has left end point \(q\); choose \(\varepsilon>0\) with
\((q-\varepsilon,q)\subseteq I_k\) and \((q,q+\varepsilon)\subseteq I_{k-1}\).
Each of these intervals contains irrationals arbitrarily close to \(q\), and
\(f_{A^*}\) takes the value determined by the parity of \(k\) on the first
and by the parity of \(k-1\) on the second.  These values differ, so
\(f_{A^*}(\beta)\) has no limit as \(\beta\to q\) and \(q\in\operatorname{Sp}(A^*)\).

\emph{The point \(1/3\).}  For every \(\varepsilon>0\) the interval
\((\tfrac13,\tfrac13+\varepsilon)\) contains \(I_k\) for every
\(k>1/\varepsilon\), hence contains irrationals at which \(f_{A^*}=1\) (even
\(k\)) and irrationals at which \(f_{A^*}=0\) (odd \(k\)).  So no limit exists
as \(\beta\to1/3\), even through \(\beta>1/3\) alone, and
\(1/3\in\operatorname{Sp}(A^*)\).

\emph{No other point of \([1/3,1)\).}  A rational \(q\in(1/3,1)\) not of the
form \(1/3+1/k\) lies in some open interval \(I_k\), on whose irrationals
\(f_{A^*}\) is constant; so \(f_{A^*}(\beta)\) converges as \(\beta\to q\) and
\(q\notin\operatorname{Sp}(A^*)\).
\end{proof}
\end{annexitem}

\begin{annexitem}{annex:p03-w03-p03-069-not-first-order}{thm:p03-w03-p03-069-not-first-order}{\(A^*\) is not first-order}
\begin{proof}
For \(m\geq1\) let \(B_{3,m}\) be the complete bipartite graph with parts
\(P\) of size \(3\) and \(Q\) of size \(m\): its edges are exactly the pairs
with one end in each part.  Suppose \(m\geq4\).  For three distinct vertices
of \(P\) the common neighbourhood is \(Q\), of size \(m\).  For three
distinct vertices of \(Q\) it is \(P\), of size \(3\).  A triple meeting both
parts has no common neighbour, since a vertex adjacent to an element of \(P\)
lies in \(Q\) and a vertex adjacent to an element of \(Q\) lies in \(P\).  So
\(\nu(B_{3,m})=m\), and by \cref{thm:p03-w03-p03-069-even-bigger},
\(B_{3,m}\entails A^*\) exactly when \(m\) is even.

Let \(\mathcal C\) be the class of finite graphs satisfying \(A^*\).  Fix a
number of rounds \(r\geq1\), let \(m=2\max(r,2)\), and compare
\(\mathcal A_r=B_{3,m}\in\mathcal C\) with
\(\mathcal B_r=B_{3,m+1}\notin\mathcal C\) in the game
\(G_r(\mathcal A_r,\mathcal B_r)\) of \cref{def:p03-ef-game}.  Fix a
bijection \(\beta\) between the two \(P\)-parts.  Duplicator answers a vertex
of a \(P\)-part by its image or preimage under \(\beta\).  She answers a vertex
of a \(Q\)-part by the vertex already paired with it, if it was chosen before,
and otherwise by a vertex of the other \(Q\)-part not chosen before; one
exists because at most \(r-1\) vertices were chosen before and each
\(Q\)-part has at least \(m\geq r\) vertices.  The resulting pairing respects
equality, because repeated choices receive repeated answers and new choices
receive new answers, and it maps each part to the corresponding part.  In a
complete bipartite graph two vertices are adjacent exactly when they lie in
different parts, so the pairing preserves and reflects adjacency.  Hence it
is a partial isomorphism and Duplicator wins.

If a first-order sentence \(\varphi\) held in exactly the finite graphs in
\(\mathcal C\), then \(\varphi\wedge\gamma\), with \(\gamma\) the
first-order sentence saying that \(\sim\) is irreflexive and symmetric, would
hold in exactly the finite \(\{\sim\}\)-structures in \(\mathcal C\).  Since
\(r\) was arbitrary, this contradicts the game criterion
\cref{thm:p03-w03-p03-038-criterion}.  So no first-order sentence defines
\(A^*\) on finite graphs.

For the consequence, suppose a first-order sentence \(\psi\) agrees with
\(A^*\) on every finite graph except those isomorphic to one of finitely many
finite graphs \(G_1,\ldots,G_m\).  Each \(G_i\), with vertices
\(g_1,\ldots,g_N\), is described up to isomorphism by the first-order sentence
\(\chi_i\) saying that there are distinct \(w_1,\ldots,w_N\), adjacent exactly
when the corresponding \(g\)'s are, such that every vertex equals some
\(w_j\).  Let \(J\) be the set of \(i\) with \(G_i\entails\psi\) and \(J^*\) the
set of \(i\) with \(G_i\entails A^*\).  The sentence
\[
  \Bigl(\psi\wedge\bigwedge_{i\in J}\neg\chi_i\Bigr)\vee\bigvee_{i\in J^*}\chi_i
\]
agrees with \(\psi\), hence with \(A^*\), on graphs isomorphic to no \(G_i\),
and on a graph isomorphic to \(G_i\) it is true exactly when \(i\in J^*\).  It
would define \(A^*\) on finite graphs, which was just excluded.
\end{proof}
\end{annexitem}

\begin{annexitem}{annex:p03-w03-p03-069-configurations}{thm:p03-w03-p03-069-configurations}{The rooted graphs behind the representation}
\begin{proof}
\emph{Item 1.}  The nonroots are \(u\) and the \(l\) vertices \(v_{ij}\), so
\(v=1+l\).  The edges not joining two roots are the three edges at each
\(v_{ij}\), so \(e=3l\).  Then \(v-e\alpha=1+l-3l\alpha=1-l(3\alpha-1)\), which
is positive exactly when \(l(3\alpha-1)<1\); it is never \(0\), since
\(l(3\alpha-1)=1\) would make \(\alpha=(l+1)/(3l)\) rational.

Suppose \(l(3\alpha-1)<1\), and let \(S'\) satisfy
\(S\subsetneq S'\subseteq V(H)\), so \(S'=S\cup X\) with \(X\) a nonempty set
of nonroots.  Count the edges of \(H|_{S'}\) not joining two roots.  If
\(u\notin X\), each \(v_{ij}\in X\) keeps only its two edges to roots, so the
subextension has \(|X|\) nonroots and \(2|X|\) edges, and
\(|X|-2|X|\alpha=|X|(1-2\alpha)>0\) because \(\alpha<1/2\).  If \(u\in X\)
together with \(j\) of the \(v_{ij}\), then there are \(1+j\) nonroots and
\(3j\) edges, and \(1+j-3j\alpha=1-j(3\alpha-1)\geq1-l(3\alpha-1)>0\) because
\(j\leq l\) and \(3\alpha-1>0\).  So every subextension is sparse and
\((S,H)\) is safe.

\emph{Item 2.}  With roots \(x_1,x_2,x_3\) and one nonroot adjacent to all
three, \(v=1\), \(e=3\) and \(v-e\alpha=1-3\alpha<0\) because \(\alpha>1/3\);
so the rooted graph is dense.  Write \(Z=\{x_1,x_2,x_3\}\) for its root
set.  The only set \(S\) with \(Z\subseteq S\subsetneq V(H)\) is \(S=Z\),
and the rooted graph \((Z,H)\) itself was just shown to be dense; so it is
rigid in the sense of \cref{def:p03-w03-p03-065-dense-sparse}.
\end{proof}
\end{annexitem}

\begin{annexitem}{annex:p03-w03-p03-069-representation}{thm:p03-w03-p03-069-representation}{Representation lemma, almost sure form}
\begin{proof}
\emph{A binomial tail bound.}  Let \(Y\) be binomial with \(N\) trials and
success probability \(q\), and \(\mu=Nq\).  We show
\[
  \Pr[Y\geq2\mu]\leq e^{-\mu/3},\qquad\Pr[Y\leq\mu/2]\leq e^{-\mu/7}.
\]
For real \(\lambda\), independence and \(1+x\leq e^x\) give
\(\mathbb E[e^{\lambda Y}]=(1+q(e^\lambda-1))^N\leq\exp(\mu(e^\lambda-1))\).
With \(\lambda=\ln2\) and Markov's inequality,
\(\Pr[Y\geq2\mu]\leq e^{-2\mu\ln2}\,\mathbb E[e^{\lambda Y}]\leq
\exp(-\mu(2\ln2-1))\), and \(2\ln2-1>0.38>1/3\).  With \(\lambda=-\ln2\),
\(\Pr[Y\leq\mu/2]=\Pr[e^{\lambda Y}\geq e^{\lambda\mu/2}]\leq
e^{\mu\ln2/2}\exp(-\mu/2)=\exp(-\mu(1-\ln2)/2)\), and
\((1-\ln2)/2>0.15>1/7\).

\emph{Global events.}  Fix an integer \(K\geq1\) with
\(\alpha>\tfrac13+\tfrac1{K+1}\), possible since \(\alpha>1/3\).  Consider
the events
\begin{enumerate}
\item[\(\mathrm E_1\):] every vertex has fewer than \(2np\) neighbours;
\item[\(\mathrm E_2\):] every two distinct vertices have at least
\((n-2)p^2/2\) and at most \(2np^2\) common neighbours;
\item[\(\mathrm E_3\):] every three distinct vertices have at most \(K\)
common neighbours.
\end{enumerate}
Each holds almost surely.  A degree is binomial with mean
\(\mu_1=(n-1)p\geq n^{1-\alpha}/2\); by the tail bound and the union bound
over \(n\) vertices, \(\Pr[\neg\mathrm E_1]\leq ne^{-\mu_1/3}\to0\).  The
number of common neighbours of two vertices is binomial with \(n-2\) trials
and success probability \(p^2\), with mean
\(\mu_2=(n-2)p^2\geq n^{1-2\alpha}/2\) for \(n\geq4\), and
\(n^{1-2\alpha}\to\infty\) because \(\alpha<1/2\); so
\(\Pr[\neg\mathrm E_2]\leq n^2(e^{-\mu_2/3}+e^{-\mu_2/7})\to0\).  Finally
\(\neg\mathrm E_3\) says that a \(K_{3,K+1}\) exists, and
\(n^{-\alpha}\ll n^{-1/3-1/(K+1)}\), so
\(\Pr[\neg\mathrm E_3]\to0\) by \cref{thm:p03-w03-p03-069-k3k-threshold}.

\emph{One root tuple.}  Let the pairs \(\{i,j\}\) with \(R(i,j)\) be
\(\pi_1,\ldots,\pi_l\), call them designated, and call the other pairs
\(\{i,j\}\subseteq\{1,\ldots,s\}\), \(i\neq j\), undesignated.  Fix distinct
vertices \(x_1,\ldots,x_s\), let \(Y=\{x_1,\ldots,x_s\}\), and condition on
the indicators of all pairs meeting \(Y\).  The following sets are then
fixed:
\begin{itemize}
\item \(N_Y\), the vertices outside \(Y\) adjacent to some \(x_i\);
\item \(U=[n]\setminus(Y\cup N_Y)\);
\item for \(i\neq j\), \(C_{ij}\), the vertices other than \(x_i,x_j\)
adjacent to both;
\item \(D\), the union of the sets \(C_{ij}\setminus Y\) over undesignated
pairs \(\{i,j\}\);
\item for \(t\leq l\) with \(\pi_t=\{i,j\}\),
\(C'_t=C_{ij}\setminus\bigl(Y\cup D\cup\bigcup_{t'\neq t}C_{\pi_{t'}}\bigr)\).
\end{itemize}
The indicators of pairs not meeting \(Y\) are independent of the conditioning
and of one another.  Let \(Q_Y\) be the event, determined by the conditioning,
that each \(x_i\) has fewer than \(2np\) neighbours, that
\((n-2)p^2/2\leq|C_{ij}|\leq2np^2\) for all \(i\neq j\), and that any three
distinct elements of \(Y\) have at most \(K\) common neighbours.  Then
\(\mathrm E_1\wedge\mathrm E_2\wedge\mathrm E_3\) implies \(Q_Y\).

Assume \(Q_Y\).  First, \(|N_Y|<2snp\), so \(|U|\geq n-s-2snp\geq n/2\) for
large \(n\), since \(np=n^{1-\alpha}\).  Second, \(|D|\leq s^2\cdot2np^2\).
Third, let \(\pi_t=\{i,j\}\).  A vertex of \(C_{ij}\) that also lies in
\(C_{\pi'}\) for another pair \(\pi'\) is adjacent to every vertex of
\(\pi_t\cup\pi'\), which contains at least three distinct roots; so it is a
common neighbour of three distinct elements of \(Y\), and there are at most
\(K\) such vertices for each \(\pi'\).  The same holds for vertices of
\(C_{ij}\cap D\), each of which lies in \(C_{\pi'}\) for an undesignated
\(\pi'\neq\pi_t\).  Fewer than \(s^2\) pairs \(\pi'\) exist, so
\[
  |C'_t|\geq\frac{(n-2)p^2}{2}-s-s^2K\geq\frac{(n-2)p^2}{4}
\]
for large \(n\).  Let \(m=\lceil(n-2)p^2/4\rceil\) and let \(C''_t\) consist of
the \(m\) least elements of \(C'_t\).  By construction the sets
\(D\), \(C''_1,\ldots,C''_l\) are pairwise disjoint subsets of \(N_Y\), and
\(U\) is disjoint from \(N_Y\).

For \(u\in U\) let \(\mathrm{Good}_u\) be the event that \(u\) has no
neighbour in \(D\) and, for each \(t\leq l\), at least one neighbour in
\(C''_t\).

\emph{\(\mathrm{Good}_u\) gives the representation.}  Let \(\pi_t=\{i,j\}\)
be designated.  A neighbour \(v\in C''_t\) of \(u\) is adjacent to \(x_i\),
\(x_j\) and \(u\), so \(R_u(x_i,x_j)\).  Let \(\{i,j\}\) be undesignated and
suppose some \(v\) is adjacent to \(x_i\), \(x_j\) and \(u\).  Then
\(v\neq x_i,x_j\), so \(v\in C_{ij}\).  If \(v\in Y\), then \(u\) is adjacent
to an element of \(Y\) and \(u\in N_Y\), contradicting \(u\in U\).  Otherwise
\(v\in C_{ij}\setminus Y\subseteq D\), contradicting \(\mathrm{Good}_u\).  So
\(\neg R_u(x_i,x_j)\).

\emph{Independence and probability.}  \(\mathrm{Good}_u\) depends only on the
indicators of the pairs \(\{u,w\}\) with \(w\in D\cup C''_1\cup\cdots\cup
C''_l\).  None of these pairs meets \(Y\), because \(u\in U\) and \(w\in
N_Y\).  For distinct \(u,u'\in U\) the two sets of pairs are disjoint,
because \(w\in N_Y\) is never in \(U\), so \(u\) is the only end of
\(\{u,w\}\) in \(U\).  Hence, given the conditioning, the events
\(\mathrm{Good}_u\), \(u\in U\), are independent, and
\[
  \Pr[\mathrm{Good}_u\mid\cdot\,]=(1-p)^{|D|}\prod_{t=1}^{l}\bigl(1-(1-p)^{m}\bigr).
\]
By Bernoulli's inequality, \((1-p)^{|D|}\geq1-p|D|\geq1-2s^2np^3
=1-2s^2n^{1-3\alpha}\geq\tfrac12\) for large \(n\), since \(\alpha>1/3\).  For
\(x\geq0\), \(1-e^{-x}\geq x-x^2/2\), because the difference vanishes at
\(0\) and has derivative \(e^{-x}-1+x\geq0\); so \(1-e^{-x}\geq x/2\) when
\(0\leq x\leq1\).  Here \(pm\leq p\bigl((n-2)p^2/4+1\bigr)\leq
n^{1-3\alpha}/4+n^{-\alpha}\leq1\) for large \(n\), and
\(pm\geq(n-2)p^3/4\geq n^{1-3\alpha}/8\) for \(n\geq4\).  Using
\(1-p\leq e^{-p}\),
\[
  1-(1-p)^m\geq1-e^{-pm}\geq\frac{pm}{2}\geq\frac{n^{1-3\alpha}}{16}.
\]
So \(\Pr[\mathrm{Good}_u\mid\cdot\,]\geq\varrho_n:=\tfrac12\,
16^{-l}\,n^{-l(3\alpha-1)}\) for every \(u\in U\), and
\[
  \Pr\bigl[\text{no }u\in U\text{ satisfies }\mathrm{Good}_u\bigm|\cdot\,\bigr]
  \leq(1-\varrho_n)^{n/2}\leq e^{-\varrho_nn/2}=e^{-cn^{\delta}},
\]
where \(c=\tfrac14\,16^{-l}\) and \(\delta=1-l(3\alpha-1)>0\).

\emph{All root tuples.}  The bounds for large \(n\) above depend only on
\(\alpha\), \(s\), \(l\) and \(K\), not on the tuple.  Let \(F\) be the event
that some distinct \(x_1,\ldots,x_s\) have no \(u\) with the required
property.  If \(F\) holds together with \(\mathrm E_1\wedge\mathrm
E_2\wedge\mathrm E_3\), then for some tuple \(Q_Y\) holds and no
\(\mathrm{Good}_u\) occurs.  Averaging the conditional bound over the
conditioning, on \(Q_Y\), gives for each tuple a probability at most
\(e^{-cn^\delta}\).  There are at most \(n^s\) tuples, so for large \(n\)
\[
  \Pr[F]\leq\Pr[\neg(\mathrm E_1\wedge\mathrm E_2\wedge\mathrm E_3)]
  +n^se^{-cn^{\delta}},
\]
where the first term tends to \(0\) and so does the second.
\end{proof}
\end{annexitem}

\begin{annexitem}{annex:p03-w03-p03-069-representation-fails}{thm:p03-w03-p03-069-representation-fails}{The representation above \(1/2\)}
\begin{proof}
\emph{Item 1.}  Fix an integer \(K'\geq1\) with \((K'+1)(2\alpha-1)>2\),
possible since \(\alpha>1/2\).  Two events hold almost surely.  First, every
vertex has fewer than \(2np\) neighbours: the degree of a vertex is binomial
with mean \(\mu_1=(n-1)p\geq n^{1-\alpha}/2\to\infty\), and the binomial tail
bound proved in the annex item for
\cref{thm:p03-w03-p03-069-representation}, with the union bound over \(n\)
vertices, bounds the failure probability by \(ne^{-\mu_1/3}\to0\).  Second,
every two distinct vertices have at most \(K'\) common neighbours: the
expected number of pairs \(\{a,b\}\) together with \(K'+1\) further vertices
all adjacent to both \(a\) and \(b\) is at most
\(n^{K'+3}p^{2(K'+1)}=n^{(K'+1)(1-2\alpha)+2}\to0\), and the first moment
bound of \cref{thm:p03-w03-p03-065-moments} applies.  Assume both events,
and let \(x_1,\ldots,x_s\) be distinct.  Let \(C\) be the set of vertices
adjacent to at least two of the \(x_i\); then \(|C|\leq s^2K'\).  Fewer than
\(s^2K'\cdot2np\) vertices are adjacent to some element of \(C\), and
\(2s^2K'np=2s^2K'n^{1-\alpha}<n\) for large \(n\), so some vertex \(u\) has
no neighbour in \(C\).  If some \(v\) were adjacent to \(x_i\), \(x_j\) and
\(u\) with \(i<j\), then \(v\in C\) and \(u\) would have a neighbour in
\(C\).  So \(\neg R_u(x_i,x_j)\) for all \(i<j\).

\emph{Item 2.}  Since \(\alpha>1/2\), \(p^2n=n^{1-2\alpha}\to0\), so
\(p^2n\leq\ln n\) for all large \(n\).  By part~2 of
\cref{thm:p03-w03-p03-065-common-neighbour}, with \(\varepsilon=1\), the
property that every two distinct vertices have a common neighbour holds almost
never.  So almost surely there are distinct \(x_1,x_2\) with no common
neighbour.  For such \(x_1,x_2\) and any vertex \(u\), a witness \(v\) for
\(R_u(x_1,x_2)\) would be a common neighbour of \(x_1\) and \(x_2\); hence
\(R_u(x_1,x_2)\) fails for every \(u\).  The sentence of the source's lemma
with \(s=2\) and \(R(1,2)\) is
\(\forall_{x_1,x_2}\exists_u\,(R_u(x_1,x_2)\leftrightarrow R(1,2))\), that
is, \(\forall_{x_1,x_2}\exists_u\,R_u(x_1,x_2)\) over distinct
\(x_1,x_2\), and it therefore holds almost never.

\emph{Consequence.}  \(I_1=(5/6,1)\), \(I_2=(2/3,5/6)\), \(I_3=(7/12,2/3)\),
\(I_4=(8/15,7/12)\) and \(I_5=(1/2,8/15)\) all lie in \((1/2,1)\).  For
\(k\leq3\), \(l<k/3\leq1\) forces \(l=0\), and item~1 applies; the
conjunction in the source's sentence then reads
\(\bigwedge_{i<j}\neg R_u(x_i,x_j)\).  For \(k\in\{4,5\}\), \(l=1<k/3\) is
admissible, and item~2 applies.
\end{proof}
\end{annexitem}

\begin{annexitem}{annex:p03-w03-p03-069-encoding}{thm:p03-w03-p03-069-encoding}{The almost sure encoding}
\begin{proof}
Throughout, \(R_u(y,w)\) is symmetric in \(y,w\), since its definition is.

\emph{Part 1, EVEN.}  Suppose \(\operatorname{EVEN}^u(\bar x)\) holds with
witnesses \(u_1,u_2\), and let \(S=N(\bar x)\).  For \(y,y'\in S\) let
\(P(y,y')\) hold when \(y\neq y'\) and
\(R_{u_1}(y,y')\vee R_{u_2}(y,y')\).  This relation on \(S\) is irreflexive
and symmetric, and by the formula every \(y\in S\) has exactly one
\(y'\in S\) with \(P(y,y')\).  By the first half of the EVEN part of
\cref{thm:p03-w03-p03-069-even-bigger}, applied to \(P\), \(|S|\) is even.

\emph{Part 1, BIGGER.}  Suppose \(\operatorname{BIGGER}^u(\bar x;\bar z)\)
holds with witnesses \(u_1,\ldots,u_4\), and let \(S=N(\bar x)\),
\(T=N(\bar z)\).  Clause (a) gives, for each \(y\in T-S\), some \(f(y)\in
S-T\) with \(R^*(y,f(y))\).  If \(f(y_1)=f(y_2)=w\) with \(y_1\neq y_2\), then
\(y_1,y_2,w\) violate clause (b).  So \(f\) is injective,
\(|T-S|\leq|S-T|\), and \(|T|\leq|S|\) as in
\cref{thm:p03-w03-p03-069-even-bigger}.

\emph{Part 2: the almost sure event.}  Let \(k\geq6\) and \(\alpha\in I_k\).
Then \(\tfrac13<\alpha<\tfrac12\), since \(I_k\subseteq(\tfrac13,\tfrac12]\)
is open, and \(3\alpha-1<3/k\).  Put \(L=\lceil k/4\rceil\).  Then
\(L\leq k/3\): for \(k=6,7,8\), \(L=2\leq k/3\); for \(k\geq9\),
\(L\leq(k+3)/4\leq k/3\).  Hence every \(l\leq L\) satisfies
\(l(3\alpha-1)<3l/k\leq1\).  Let \(\mathrm E\) be the event that
\(\nu(G)\leq k\) and that, for every \(s\) with \(2\leq s\leq2k\) and every
symmetric irreflexive \(R\) on \(\{1,\ldots,s\}\) holding for at most \(L\)
unordered pairs, every distinct \(w_1,\ldots,w_s\) have a vertex \(u\) with
\(R_u(w_i,w_j)\leftrightarrow R(i,j)\) for all \(i<j\).  This is a
conjunction of finitely many events, each holding almost surely by
\cref{thm:p03-w03-p03-069-max-common} and
\cref{thm:p03-w03-p03-069-representation}, so \(\mathrm E\) holds almost
surely.  Assume \(\mathrm E\) and \(n\geq3\) for the rest of the proof.

\emph{Part 2, EVEN.}  Let \(\bar x\) be distinct and \(S=N(\bar x)\), so
\(|S|=s\leq k\), and suppose \(s\) is even.  If \(s=0\), the formula holds
with any \(u_1,u_2\), as there is no \(y\) to check.  If \(s\geq2\), list
\(S\) as \(w_1,\ldots,w_s\) and let \(M\) be the perfect matching
\(\{w_1,w_2\},\{w_3,w_4\},\ldots\) with \(s/2\) pairs.  Let \(M_1\) be its
first \(\lceil s/4\rceil\) pairs and \(M_2\) the remaining
\(\lfloor s/4\rfloor\); each has at most \(\lceil k/4\rceil=L\) pairs.  By
\(\mathrm E\), with \(R\) the relation given by \(M_1\) on the indices, there
is \(u_1\) such that for \(i<j\), \(R_{u_1}(w_i,w_j)\) holds exactly when
\(\{w_i,w_j\}\in M_1\); likewise there is \(u_2\) for \(M_2\).  Then for
distinct \(y,y'\in S\), \(R_{u_1}(y,y')\vee R_{u_2}(y,y')\) holds exactly
when \(\{y,y'\}\in M\), so every \(y\in S\) has exactly one \(y'\in S\),
\(y'\neq y\), with that property, namely its partner in \(M\).  So
\(\operatorname{EVEN}^u(\bar x)\) holds.  With Part 1,
\(\operatorname{EVEN}^u(\bar x)\) holds exactly when \(|N(\bar x)|\) is even.

\emph{Part 2, BIGGER.}  Let \(\bar x,\bar z\) be distinct triples,
\(S=N(\bar x)\), \(T=N(\bar z)\), and suppose \(|T|\leq|S|\).  Then
\(|T-S|\leq|S-T|\), and both have at most \(k\) elements.  If \(T-S\) is
empty, clauses (a) and (b) hold for any \(u_1,\ldots,u_4\).  Otherwise take
an injection \(f\colon T-S\to S-T\), list the at most \(2k\) distinct
vertices of \((T-S)\cup(S-T)\) as \(w_1,\ldots,w_{s}\), \(s\geq2\), and split
the \(|T-S|\leq k\) pairs \(\{y,f(y)\}\) into four parts with at most
\(\lceil|T-S|/4\rceil\leq L\) pairs each.  By \(\mathrm E\) there are
\(u_1,\ldots,u_4\) such that \(R_{u_h}\) holds on a pair of distinct
\(w\)'s exactly when the pair is in the \(h\)-th part.  So for distinct
\(y,w\) in \((T-S)\cup(S-T)\), \(R^*(y,w)\) holds exactly when
\(\{y,w\}=\{y',f(y')\}\) for some \(y'\in T-S\).  Clause (a) holds with
\(w=f(y)\).  For clause (b), if \(y_1\neq y_2\) in \(T-S\), \(w\in S-T\),
and \(R^*(y_1,w)\), \(R^*(y_2,w)\), then, since \(T-S\) and \(S-T\) are
disjoint, \(w=f(y_1)=f(y_2)\), contradicting injectivity.  So
\(\operatorname{BIGGER}^u(\bar x;\bar z)\) holds.  With Part 1, it holds
exactly when \(|N(\bar z)|\leq|N(\bar x)|\).

\emph{Part 2, the sentences.}  Under \(\mathrm E\), \(A\) holds exactly when
some distinct \(\bar x\) has \(|N(\bar x)|\) even and
\(|N(\bar z)|\leq|N(\bar x)|\) for all distinct \(\bar z\).  By
\cref{thm:p03-w03-p03-069-even-bigger} this is exactly \(A^*\).  So
\(\Pr[A\leftrightarrow A^*]\geq\Pr[\mathrm E\wedge n\geq3]\to1\).
\end{proof}
\end{annexitem}

\begin{annexitem}{annex:p03-w03-p03-069-infinite-spectrum}{thm:p03-w03-p03-069-infinite-spectrum}{A first-order sentence with infinite spectrum}
\begin{proof}
Let \(k\geq6\) and \(\alpha\in I_k\) irrational.  Since
\(|\Pr[A]-\Pr[A^*]|\leq\Pr[A\not\leftrightarrow A^*]\), which tends to \(0\)
by \cref{thm:p03-w03-p03-069-encoding}, \(f_A(\alpha)=f_{A^*}(\alpha)\), and
this is \(1\) for even \(k\) and \(0\) for odd \(k\) by
\cref{thm:p03-w03-p03-069-astar-spectrum}.

For \(k\geq7\) let \(q=\tfrac13+\tfrac1k\).  Points just below \(q\) lie in
\(I_k\) and points just above lie in \(I_{k-1}\), with \(k-1\geq6\); so near
\(q\) the function \(f_A\) equals \(1\) on the irrationals of one side and
\(0\) on the irrationals of the other, it has no limit at \(q\), and
\(q\in\operatorname{Sp}(A)\).  Every interval \((\tfrac13,\tfrac13+\varepsilon)\)
contains \(I_k\) for all large \(k\), so contains irrationals where
\(f_A=1\) and irrationals where \(f_A=0\); hence \(1/3\in\operatorname{Sp}(A)\).
A rational \(q\in(1/3,1/2)\) not of the form \(1/3+1/k\) lies in some open
\(I_k\) with \(k\geq6\), on whose irrationals \(f_A\) is constant, so
\(q\notin\operatorname{Sp}(A)\).  The points \(1/3+1/k\), \(k\geq7\), are
infinitely many and distinct, so \(\operatorname{Sp}(A)\) is infinite.
\end{proof}
\end{annexitem}
